\documentclass[11pt]{article}

\usepackage[a4paper,margin=1in]{geometry}
\usepackage{amsmath,amssymb,amsthm}
\usepackage{enumitem}
\usepackage{xcolor}
\usepackage{fancyhdr}
\usepackage{needspace}

\definecolor{courseblue}{RGB}{0,0,255}

\setlength{\parindent}{0pt}
\setlength{\parskip}{0.35em}
\setlength{\headheight}{14pt}

\pagestyle{fancy}
\fancyhf{}
\fancyfoot[C]{\thepage}
\renewcommand{\headrulewidth}{0pt}

\newcounter{problem}
\newenvironment{problem}{%
    \refstepcounter{problem}%
    \Needspace{7\baselineskip}%
    \par\vspace{0.7em}%
    \noindent\textcolor{courseblue}{\rule{\textwidth}{0.4pt}}%
    \par\vspace{0.5em}%
    \begin{list}{}{%
        \setlength{\leftmargin}{2.3em}%
        \setlength{\labelwidth}{1.8em}%
        \setlength{\labelsep}{0.5em}%
        \setlength{\topsep}{0pt}%
        \setlength{\partopsep}{0pt}%
    }%
    \item[{[\theproblem]}]
}{%
    \end{list}
    \par\vspace{0.5em}
}

\newenvironment{parts}{%
    \begin{enumerate}[label={[\alph*]},leftmargin=2.2em,itemsep=0.65em,topsep=0.65em]
}{%
    \end{enumerate}
}

% Type each solution between \begin{answer} and \end{answer}.
\newenvironment{answer}{%
    \Needspace{4\baselineskip}%
    \par\vspace{0.35em}
    \noindent\textbf{Answer:}\par
}{%
    \par\vspace{0.25em}
}

\begin{document}

\begin{center}
    {\color{courseblue}\Large\textbf{AMATH/PMATH 331: Applied Real Analysis}}
    \hfill
    {\color{courseblue}\Large\textbf{Fall 2026}}

    \vspace{1.2em}
    {\color{courseblue}\Large
    Assignment 1 Submission; due Thursday, September 24, 2026}

    \vspace{1.2em}
    
    \textbf{Name:} Vlad Paun

    \textbf{Student number:} 21291453
\end{center}

\begin{problem}
Let $E \subseteq \mathbb{R}$ be nonempty. We say that $t \in \mathbb{R}$ is a
\emph{lower bound} of $E$ if $t \leq x$ for all $x \in E$. We say that $E$ is
\emph{bounded below} if there exists a lower bound of $E$. We say that
$t \in \mathbb{R}$ is a \emph{greatest lower bound}, or \emph{infimum}, of $E$
if it is a lower bound of $E$ and if $t \geq s$ for any lower bound $s$ of $E$.
If $E$ has an infimum, it is denoted $\inf E$.

\begin{parts}
    \Needspace{8\baselineskip}\item Let $E$ be a nonempty subset of $\mathbb{R}$ which is bounded below.
    Find a way to use the completeness axiom to prove that $\inf E$ exists.
    \emph{Hint:} Consider the subset $-E = \{-x : x \in E\}$.

    \begin{answer}
        \begin{proof}
            Since $E$ is nonempty and bounded below, there exists
            $\alpha \in \mathbb{R}$ such that
            \[
                \forall x \in E:\ \alpha \leq x
            \]

            Consequently, 
            \[
                \forall x \in E:\ -\alpha \geq -x
            \].

            Define
            \[
                -E=\{-x:x\in E\}.
            \]
            
            Then, $-\alpha$ is an upper bound of $-E$.


            By the completeness axiom, $\sup(-E)$ exists. Let
            \[
                \beta=\sup(-E).
            \]

            It follows that $\forall x \in -E:\ \beta \geq x$. And therefore, $\forall x \in E:\ -\beta \leq x$. So $-\beta$ is a lower bound of $E$.

            Next, take any other lower bound of $E$, call it $s$. From the same logic as above, it follows that $-s$ is an upper bound of $-E$. Since $\beta$ is a supremum of $-E$, it follows that $\beta \leq -s$. 
            
            So, $-\beta \geq s$. Therefore, since $-\beta$ is a lower bound of $E$ and it is at least as large as any other lower bound, it follows that $-\beta = \inf E$ 
        \end{proof}
    \end{answer}

    \Needspace{8\baselineskip}\item Let $\varnothing \neq A \subseteq B \subseteq \mathbb{R}$. Suppose
    $B$ is bounded below (so $\inf B$ exists). Prove that $A$ is bounded below
    (so $\inf A$ exists) and that $\inf A \geq \inf B$.

    \begin{answer}
        \begin{proof}
            If $B$ is bounded below, then $\exists \alpha:\ \forall x \in B:\ \alpha \leq x$.

            Since $A \subseteq B$, we know that $\forall y\in A:\ y\in B$.

            Consequently,
            \[
            \forall y \in A: \alpha \leq y
            \]

            Therefore, $\alpha$ is a also a lower bound of $A$. And from $1[a]$, we can conclude that $\inf A$ also exists.
            
            Now, since $A$ is a subset of $B$ and $\inf B$ is a lower bound of $B$, then $\inf B$ must also be a lower bound of $A$. Since $\inf A$ must be greater then or equal to any other lower bounds of $A$, we have that
            \[
            \inf A \geq \inf B
            \]
        \end{proof}
    \end{answer}

    \Needspace{8\baselineskip}\item Let $A,B$ be nonempty subsets of $\mathbb{R}$, both bounded below
    (so both $\inf A$ and $\inf B$ exist). Define
    \[
        A+B = \{a+b : a \in A,\ b \in B\}.
    \]
    Prove that $A+B$ is bounded below and that
    $\inf(A+B)=\inf A+\inf B$.

    \begin{answer}
        \begin{proof}
            \begin{align*}
                \forall a\in A&:\ \inf A \leq a \\
                \forall b\in B&:\ \inf B \leq b \\
                \therefore \forall a\in A,b\in B&:\ \inf A + \inf B \leq a+b\\
                \therefore \forall x\in A+B&:\ \inf A + \inf B \leq x
            \end{align*}

            So $\inf A + \inf B$ is a lower bound of $A+B$. Consequently, $\inf(A+B)$ exists.

            Now, we know that
            \begin{align*}
                \forall \varepsilon > 0&:\ \exists \alpha\in A:\ \inf A + \frac{\varepsilon}{2}  > \alpha \\
                \forall \varepsilon > 0&:\ \exists \beta\in B:\ \inf B + \frac{\varepsilon}{2}  > \beta \\
            \end{align*}

            Therefore,
            \begin{align*}
                \forall \varepsilon > 0&:\ \exists \alpha\in A, \beta\in B:\ \inf A + \inf B + \varepsilon  > \alpha + \beta \\
            \end{align*}
            This means that no number greater then $\inf A + \inf B$ can be a lower bound of $A+B$. Therefore, 
            \[
            \inf(A+B) = \inf A + \inf B
            \]
        \end{proof}
    \end{answer}
\end{parts}
\end{problem}

\begin{problem}
Let $x,y$ be two real numbers such that $y-x>1$. (In particular, $x<y$.)

\begin{parts}
    \Needspace{8\baselineskip}\item If $x$ is positive, prove that there exists an integer $k \in \mathbb{Z}$
    such that $x<k<y$.

    \begin{answer}
        \begin{proof}
            By the Archimedean property, the set $E=\{n\in\mathbb{N}:\ n>x\}\subset \mathbb{N}\subset \mathbb{Z}$ is not empty.

            $x$ is a lower bound of $E$, so by the completeness axiom, $\inf E$ exists. And since $E\subset\mathbb{Z}$, $\inf E\in E$. Let $\inf E = k$.

            Then $x<k$. Now, assume for the sake of contradiction that $x<k-1$. Since $k-1\in\mathbb{Z}$, then $k-1\in E$. But $k-1 < k$ and $k = \inf E$, which is a contradiction.

            Therefore, $k-1 \leq x$. $k \leq x+1 < y$.

            So, $x < k < y$.
            
        \end{proof}
    \end{answer}

    \Needspace{8\baselineskip}\item If $x$ is negative and $y$ is positive, prove that there exists an
    integer $k \in \mathbb{Z}$ such that $x<k<y$.

    \begin{answer}
        \begin{proof}
            Take $k=0\in\mathbb{Z}$.
        \end{proof}
    \end{answer}

    \Needspace{8\baselineskip}\item If $y$ is negative, prove that there exists an integer $k \in \mathbb{Z}$
    such that $x<k<y$.

    \begin{answer}
        \begin{proof}
            Since $y<0$, we have $-y>0$. Moreover,
            \[
                (-x) - (-y) = y-x > 1 
            \]
            So applying $2[a]$ to $-x$ \& $-y$, we get that $\exists n\in\mathbb{Z}:\ -y < n < -x$.
            \[
                x < -n < y
            \]
            Taking $k=-n\in\mathbb{Z}$ proves the result.
        \end{proof}
    \end{answer}
\end{parts}
\end{problem}

\begin{problem}
The parts of this problem are \emph{not related}.

\begin{parts}
    \Needspace{8\baselineskip}\item Let $x \in \mathbb{R}$. Consider the set
    $E=\{q \in \mathbb{Q}:q<x\}\subset \mathbb{Q}\subset \mathbb{R}$.
    Show that $E$ is nonempty and bounded above, so $\sup E$ exists. Show also
    that $\sup E=x$.

    \begin{answer}
    \end{answer}

    \Needspace{8\baselineskip}\item Let $x,y \in \mathbb{R}$ with $x<y$. Prove that there exists an
    \emph{irrational} (not rational) number $z$ such that $x<z<y$.
    \emph{Hint:} You may use the fact that $\sqrt{2}$ is irrational.

    \begin{answer}
    \end{answer}
\end{parts}
\end{problem}

\begin{problem}
Let $C>0$ be a fixed positive number.

\begin{parts}
    \Needspace{8\baselineskip}\item Let $x \geq 0$, and suppose that for any $\varepsilon>0$, we have
    $x<C\varepsilon$. Show that for any $\eta>0$, we have $x<\eta$.

    \begin{answer}
    \end{answer}

    \Needspace{8\baselineskip}\item In the situation of part [a], show that we must have $x=0$.
    \emph{Hint:} Find a way to use the Archimedean property.

    \begin{answer}
    \end{answer}
\end{parts}
\end{problem}

\begin{problem}
COMING SOON

\begin{answer}
\end{answer}
\end{problem}

\begin{problem}
COMING SOON

\begin{answer}
\end{answer}
\end{problem}

\begin{problem}
COMING SOON

\begin{answer}
\end{answer}
\end{problem}

\begin{problem}
COMING SOON

\begin{answer}
\end{answer}
\end{problem}

\begin{problem}
COMING SOON

\begin{answer}
\end{answer}
\end{problem}

\begin{problem}
COMING SOON

\begin{answer}
\end{answer}
\end{problem}

\end{document}
